\documentclass{article}
\usepackage{amsmath}
\usepackage{amsthm}

\newtheorem{theorem}{Problem}

\begin{document}

\begin{theorem}
Ten treeks weigh as much as three squigs and one goolee. Two treeks and one goolee are equal in weight to one squig. The combined weight of how many treeks equals the weight of one squig?
\end{theorem}

\begin{proof}
Step 0: So first let's set up some variables. Let's let $t$ be the weight of one treek, $s$ be the weight of one squig, and $g$ be the weight of one goolee.
Step 1: Now we need to write some equations using the given information.
Step 2: The first sentence says that $10t$ is equal to $3s$ plus $g$.
Step 3: Or $10t=3s+g$
Step 4: Ok and the second sentence says that $2t$ plus $g$ is equal to $s$.
Step 5: Ok so we have the system of equations $10t=3s+g$ and $2t+g=s$
Step 6: Let's eliminate $g$ by subtracting the equations.
Step 7: No, let's rather add the equations and subtract $g=g$.
Step 8: That's a good idea. So the sum of the equations is $10t+2t+g=3s+g+s$, which simplifies to $12t=4s$.
Step 9: That means that $4t=s$.
Step 10: Yes and that's our answer because we were asked to find how many treeks are equal to one squig.
Step 11: Right, four treeks are equal to one squig.
\end{proof}

\end{document}
